\section{Fundamentação Teórica}

Este capítulo estabelece o arcabouço teórico necessário para a otimização proposta, situando o estudo na interseção entre Processamento de Linguagem Natural (NLP) \citep{jurafsky2024}, Recuperação de Informação (IR) \citep{manning2008} e Estatística Experimental \citep{montgomery2017}. A exposição segue uma estrutura lógico-construtiva: inicia-se com a representação matemática do texto em espaços vetoriais, avançando para a modelagem estocástica de sequências por meio da arquitetura \textit{Transformer}. Em seguida, detalha-se a engenharia de sistemas RAG, abrangendo estratégias de segmentação e algoritmos de busca híbrida. Por fim, definem-se os protocolos de avaliação automatizada e detalha-se o referencial analítico, fundamentado no delineamento Fatorial Completo e no método não-paramétrico Tranformação de Postos Alinhados (\textit{Aligned Rank Transform} - ART) \citep{wobbrock2011art}, que viabilizará a investigação quantitativa dos hiperparâmetros.
